\SourceLeaf{172}schnittszahl der jährlich neu beitretenden Mitglieder ausdrückt. In der Denkschrift von 1845 habe ich dafür \(2\tfrac{1}{2}\) angenommen, ohne zu verschweigen, dass dieses Element ein sehr ungewisses ist: alles wohl erwogen, habe ich dieselbe Ziffer auch diesmal beibehalten zu müssen geglaubt, und so haben sich die in der Bilanzrechnung angesetzten Positionen (8) und (3) ergeben.

Die Position (4) ist die capitalisirte jährliche Ausgabe von 900 Thalern, wovon der bewegliche Theil der Witwenpension zu bestreiten ist. Diese Summe wird unter alle berechtigten Witwen zu gleichen Theilen vertheilt, wenn deren Anzahl 18 oder mehr beträgt; ist die Anzahl kleiner, so erhält jede 50 Thaler. Es erhellet hieraus, dass im letztern Falle (der auch in diesem Augenblick Statt findet) die Kasse eine Ersparniss macht, welche in der Bilanz nicht mit berechnet ist, und für etwas längere Zeit auch gar nicht im Voraus berechnet werden kann: jedenfalls aber ist diess nur ein vorübergehender Vortheil, welcher in späterer Zeit, wenn die Folgen der jetzigen grossen Ausdehnung der Genossenschaft sich erst entwickelt haben werden, selten oder vielleicht niemals wieder vorkommen wird.

Als Durchschnittswerthe der Nebenausgaben habe ich angenommen
\begin{center}
\small\setlength{\tabcolsep}{2pt}
\begin{tabular}{r r l r l r l p{196pt}}
&28&Thl.&4&Ggr&—&Pf.&für Processkosten, soweit sie nicht erstattet,\\
&61&,,&8&,,&9&,,&für Baukosten\\
&133&,,&9&,,&5&,,&für Rechnungsführung und Copialien\\
&136&,,&6&,,&3&,,&für Verluste\\
Zusammen&359&Thl.&4&Ggr.&5&Pf.&
\end{tabular}
\end{center}
\noindent\hspace*{\parindent}Diese Ansätze gründen sich auf die Erfahrungen der letzten 21 oder 20 Jahre. Bei der Rechnung von 1845 hatten nur Erfahrungen von 15 oder 14 Jahren zum Grunde gelegt werden können, welche die Totalsumme 323 Thl. 13 Ggr. ergeben hatten. Die erstere Zahl, capitalisirt, erbringt 10262 Thl. 10 Ggr. Dieser Betrag unter Zufügung von 2644 Thl. 4 Ggr. (als mittlern Werthe des unproductiven Vermögenstheils nach 17jährigem Durchschnitt) bildet die Position (5). Für die letztere Zahl war übrigens in der Rechnung von 1845 nach 11jährigem Durchschnitt der nahe gleiche Werth 2635 Thl. 18 Ggr. 4 Pf. angenommen worden.

Die Position (6) entsteht aus der Capitalisirung der jährlichen Einnahme aus dem Pachtzins der Universitäts-Apotheke (1000 Thaler).

Die Position (9) bedarf auch noch einiger Erläuterungen. In der Jahresrechnung für 1850—1851 ist das Geldvermögen der Kasse für den 1. Julius 1851 zu 125076 Thl. 1 Ggr. 11 Pf. angesetzt, wovon die verzinslichen Capitale 123372 Thl. 2 Ggr. 7 Pf. ausmachen; das übrige besteht in dem baaren Geldvorrath, den Rückständen, und einem dem Universitätsapotheker bewilligten unverzinslichen Vorschuss. Die Capitale sind etwa zur Hälfte bei Privatschuldnern hypothekarisch, die übrigen in unkündbaren Staatspapieren angelegt, und diese letztern sind in der neuesten Jahresrechnung (eben so wie schon in mehrern vorhergehenden) schlechthin nach dem Nominalwerthe in Ansatz gebracht. Im Jahre 1845 waren hingegen diese Ansätze nach den Ankaufpreisen gemacht, und ich habe in meiner damaligen Rechnung dieselben ungeändert beibehalten, weil damals die Schwankungen in dem Werthe der Staatspapiere viel geringer waren, als seit den letzten 3 bis 4 Jahren. Jetzt, wo ein beträchtlicher Theil der zum Vermögen der Witwenkasse gehörenden Staatspapiere so sehr tief unter dem Nennwerthe steht, halte ich für nothwendig, in der Bilanzrechnung die Papiere nach dem zeitigen wirklichen Werthe, wie sie sich realisiren lassen, zu evaluiren. Ich habe dazu die Börsencourse in Frankfurt und Hannover vom 1. October angewandt, weil doch ein bestimmtes Datum gewählt werden musste. Es sind die folgenden
\begin{center}
\footnotesize\setlength{\tabcolsep}{2pt}
\begin{tabular}{l r|l r}
Oesterreichische \(4\tfrac{1}{2}\) Metalliques . .&\(67\tfrac{1}{8}\)&Hannoversche 5 Proc. . . .&\(104\tfrac{1}{2}\)\\
— \(4\tfrac{1}{2}\) bei Goll . . . .&\(69\tfrac{3}{4}\)&— 4 Proc. . . .&102\\
— \(4\tfrac{1}{2}\) bei Bethmann&\(71\tfrac{3}{4}\)&— \(3\tfrac{1}{2}\) Courant&\(98\tfrac{1}{4}\)\\
Badensche \(3\tfrac{1}{2}\) . . . . . . . . . . . . . .&\(87\tfrac{3}{4}\)&— \(3\tfrac{1}{2}\) Gold . .&\(97\tfrac{3}{4}\)
\end{tabular}
\end{center}
